% Volume I: The Epistemology of Suspicion
% Part II. Events, Records, and Reconstructions
% Chapter 8: Recording as Selection
% Spine: proof-spines/ch008-spine.md

\chapter{Recording as Selection}
\label{ch:ch008}


No recording system captures everything. A camera selects a field of view, frame rate, spectral range, resolution, exposure, compression scheme, storage duration, and metadata policy; a form selects prompts, fields, and categories. The recorded object is therefore not simply the event itself but
\[
R=S_\theta(X),
\]
where \(\theta\) is a selection regime that determines what becomes recordable at all.

This chapter introduces \textbf{selection geometry}: the pattern of invisibilities produced when a recording apparatus is optimized for one class of events rather than another. Surveillance designed to identify faces may fail to preserve gait, peripheral interaction, audio context, or environmental causation. The important question is thus not only whether a record was altered after capture, but what the architecture of capture rendered visible or invisible from the start.
